\begingroup
\small
\setlength{\tabcolsep}{5pt}
\begin{longtable}{@{}p{.04\linewidth}>{\raggedright\arraybackslash}p{.20\linewidth}>{\raggedright\arraybackslash}p{.68\linewidth}@{}}
\caption{Retrieved studies used to develop the indicator catalog}\label{tab:retrieved-studies}\\
\toprule No. & Study & Title and source \\ \midrule
\endfirsthead
\toprule No. & Study & Title and source \\ \midrule
\endhead
\bottomrule
\endfoot
1 & Goverde (2007) & Railway timetable stability analysis using max-plus system theory. \href{https://doi.org/10.1016/j.trb.2006.02.003}{Source} \\
2 & Goverde (2010) & A delay propagation algorithm for large-scale railway traffic networks. \href{https://doi.org/10.1016/j.trc.2010.01.002}{Source} \\
3 & Büker and Seybold (2012) & Stochastic modelling of delay propagation in large networks. \href{https://doi.org/10.1016/j.jrtpm.2012.10.001}{Source} \\
4 & Harrod, Cerreto and Nielsen (2019) & A closed form railway line delay propagation model. \href{https://doi.org/10.1016/j.trc.2019.02.022}{Source} \\
5 & Andersson, Peterson and Törnquist Krasemann (2013; online 2014) & Quantifying railway timetable robustness in critical points. \href{https://doi.org/10.1016/j.jrtpm.2013.12.002}{Source} \\
6 & Ma et al. (2021) & Scheduling-measure dependent modelling of delay propagation on a single railway line: A matrix transformation approach. \href{https://doi.org/10.1049/itr2.12112}{Source} \\
7 & Meng et al. (2019) & Increasing Robustness by Reallocating the Margins in the Timetable. \href{https://doi.org/10.1155/2019/1382394}{Source} \\
8 & Wu et al. (2023) & Cascading Delays for the High-Speed Rail Network Under Different Emergencies: A Double Layer Network Approach. \href{https://doi.org/10.1109/JAS.2022.105530}{Source} \\
9 & Lian et al. (2024) & Bi-Directional Delay Propagation Analysis and Modeling for High-Speed Railway Networks Under Disturbance. \href{https://doi.org/10.1109/TCSS.2023.3322747}{Source} \\
10 & Tessitore et al. (2025) & On the fragility of a train timetable. \href{https://doi.org/10.1016/j.omega.2025.103341}{Source} \\
11 & Corman et al. (2014) & Evaluating Disturbance Robustness of Railway Schedules. \href{https://doi.org/10.1080/15472450.2013.801714}{Source} \\
12 & Zhang et al. (2024) & Propagation-Based Train Rescheduling under Recoverable Delay Disturbances. \href{https://doi.org/10.1155/2024/9961840}{Source} \\
13 & Huang et al. (2026; online 2025) & Solving the railway timetable rescheduling problem with graph neural networks. \href{https://doi.org/10.1007/s40534-025-00383-7}{Source} \\
14 & Acuña-Agost et al. (2011) & SAPI: Statistical Analysis of Propagation of Incidents. A new approach for rescheduling trains after disruptions. \href{https://doi.org/10.1016/j.ejor.2011.05.047}{Source} \\
15 & Tessitore et al. (2025) & An iterative fragility-based approach to timetable design. \href{https://doi.org/10.1016/j.jrtpm.2025.100554}{Source} \\
16 & Hua et al. (2024) & Improving timetable resilience for high-speed railways through time supplement allocation. \href{https://doi.org/10.11860/j.issn.1673-0291.20230073}{Source} \\
17 & Corman and Kecman (2018) & Stochastic prediction of train delays in real-time using Bayesian networks. \href{https://doi.org/10.1016/j.trc.2018.08.003}{Source} \\
18 & Ulak et al. (2020) & Analyzing network-wide patterns of rail transit delays using Bayesian network learning. \href{https://doi.org/10.1016/j.trc.2020.102749}{Source} \\
19 & Dekker et al. (2022) & Modelling railway delay propagation as diffusion-like spreading. \href{https://doi.org/10.1140/epjds/s13688-022-00359-1}{Source} \\
20 & Ge et al. (2024; online 2023) & Modelling the cascading effects of train delay patterns and inter-train control actions with Bayesian networks. \href{https://doi.org/10.1080/23248378.2023.2194304}{Source} \\
21 & Daniotti et al. (2024) & Systemic risk approach to mitigate delay cascading in railway networks. \href{https://doi.org/10.1038/s44333-024-00012-6}{Source} \\
22 & Xu et al. (2026) & Multivariate event hypergraph diffusion model for train delay prediction. \href{https://doi.org/10.1016/j.trc.2025.105390}{Source} \\
\end{longtable}
\endgroup
